% §785's `align_peek` reads the token that starts an alignment entry with
% `get_x_token`:
%
%     procedure align_peek;
%     label restart;
%     begin restart: align_state:=1000000; @<Get the next non-blank non-call token@>;
%
%     @<Get the next non-blank non-call token@>=
%     repeat get_x_token;
%     until cur_cmd<>spacer
%
% so in TeX82 the entry's first token is expanded *before* §788's `init_col`
% backs it up and puts the u-template in front of it. e-TeX's `\protected`
% takes that back: pdfTeX leaves a protected macro alone here and expands it
% after the template has run.
%
% It matters because the u-template is where the entry's mode is decided.
% LaTeX's `array` uses `$\displaystyle#$`, and `\{` is
%
%     \{=\protected macro:
%     ->\ifmmode \lbrace \else \textbraceleft \fi .
%
% so expanding it early answers `\ifmmode` in the alignment's internal vertical
% mode and takes the text branch. `\textbraceleft` is `\char123`, which §1151
% turns into a *math* character through `\mathcode123 = "007B` — family 0,
% character 123, an en dash in OT1 rather than the brace from family 2.
%
% There is no `etex.ch` here, so the rule was measured rather than read: the
% same probe at the head of an entry, once protected and once not.
%
%     \def\pr{...}            pdfTeX expands it early -- reports the outer mode
%     \protected\def\pp{...}  pdfTeX leaves it       -- reports math
%
% The unprotected one reporting the outer mode in pdfTeX too is what says this
% is `\protected`'s doing, and not a difference in when the template arrives.
\catcode`\{=1 \catcode`\}=2 \catcode`\#=6 \catcode`\$=3 \catcode`\&=4
\font\r=cmr10 \r \hsize=300pt
\font\mi=cmmi10 \font\sy=cmsy10 \font\ex=cmex10
\textfont0=\r \scriptfont0=\r \scriptscriptfont0=\r
\textfont1=\mi \scriptfont1=\mi \scriptscriptfont1=\mi
\textfont2=\sy \scriptfont2=\sy \scriptscriptfont2=\sy
\textfont3=\ex \scriptfont3=\ex \scriptscriptfont3=\ex

\protected\def\pp{\ifmmode\message{<prot:math>}\else\message{<prot:outer>}\fi}

% An alignment whose template opens math: the protected macro waits for the `$`.
$$\halign{$\displaystyle#$&#\cr
\pp x&y\cr}$$

% The *unprotected* counterpart is deliberately not here. It is still wrong: we
% report math where pdfTeX reports the outer mode, because §791 and §785 peek
% *before* `init_span`/`init_row` push the entry's nest and we peek after. That
% is an ordering question rather than a `\protected` one, so it is tracked on
% its own rather than smuggled into this file, which would fail as a whole and
% say nothing about what does work.

% The very first entry of an alignment reaches `align_peek` by a different
% route than a later row does — `\halign`'s own start rather than after a
% `\cr` — so it is worth both.
$$\halign{$\displaystyle#$\cr
\pp a\cr}$$

% After `\noalign`, which is the third way in: `align_peek` runs again when the
% no-align group closes, and that is the path LaTeX's `\hline` uses.
$$\halign{$\displaystyle#$\cr
\pp b\cr\noalign{\kern1pt}\pp c\cr}$$

% A protected macro that is *not* first in the entry was never peeked at, so it
% expands inside the template either way — the control case that a fix which
% simply stopped expanding here would still pass.
$$\halign{$\displaystyle#$\cr
d\pp\cr}$$

% A template that does *not* open math: the protected macro still waits, and
% still reports whatever the template left it in. This separates "waits for the
% template" from "always says math".
$$\halign{#\cr
\pp e\cr}$$

% `\protected` on a macro with arguments — `MacroDelim` rather than `Macro`
% internally, and both have to be covered.
\protected\def\pq#1{\ifmmode\message{<arg:math:#1>}\else\message{<arg:outer:#1>}\fi}
$$\halign{$\displaystyle#$\cr
\pq f\cr}$$

% An expandable *primitive* is still expanded at the peek — `\protected` cannot
% be applied to one, and a fix that stopped expanding everything here would
% change this.
$$\halign{$\displaystyle#$\cr
\number\the\numexpr 6*7\relax\cr}$$
\end
